\scanpage{428}
\mspage{340}

à l'exception de $a'\Rightarrow a$ qui est un peu plus simple dans le cas actuel. On va la dégager sous forme de lemme un peu plus \unsure{général|générale} : (NB la dém. de $a'\Rightarrow a$ dans 4.13.15 avait été omise comme ``triviale'' ! cf. \unsure{pg|p.} 337 \unsure{\ldots})

\uline{Lemme} 4.13.19 : Soient $\unsure{M|\mathcal M}$ une catégorie accessible, $\underline S$ une sous-catégorie accessible de $M$, $\pi$ un cardinal utile pour $M$ et tel que $\underline S$ soit $\pi$-accessible, i.e. $\underline S$ est l'image essentielle de $\mathrm{Ind}_\pi(\underline S(\pi))$ par le foncteur pl. fidèle canonique $\mathrm{Ind}_\pi(\underline S(\pi))\longrightarrow M$, où $\underline S(\pi)=\underline S\cap M_\pi$.

On suppose qu'il existe un ordinal limite $\alpha$, tel que $\mathrm{card}\,\alpha\leq\pi$ et tel que $M_\pi$ et $\underline S\deleted{\unsure{{}_{\pi}|(\pi)}}(\pi)$ $(=S_\pi)$ soient stables par $\varinjlim$ de type $I_\alpha$, \deleted{Alors} et que le foncteur d'inclusion $\underline S(\pi)\longrightarrow M\deleted{\unsure{(\pi)}}_\pi$ y commute. Soit $\unsure{\Phi|\mathscr I|\mathscr D}$ un ensemble de catégories \deleted{de \unsure{cardinaux|cardinal}} telles que $\forall I\in\Phi$, on ait $\unsure{\mathrm{card}\,\mathrm{red}\,I|\mathrm{card}\,\mathrm{rel}\,I|\mathrm{card}\,\text{réel}\,I}\leq\pi$, et supposons que $M\deleted{\unsure{(\pi)}}_\pi$ \deleted{soit} \deleted{et $\underline S(\pi)$ soient stables} soit stable par $\varinjlim$ de type $\Phi$, et que $\underline S(\pi)$ soit stable dans $M_\pi$ par lesdites limites (ou ce qui revient au $=$, qu'il soit stable en lui-même pour ces $\varinjlim$\unsure{\textsuperscript{t}|,}, et que le foncteur d'inclusion y commute). Alors $M$ et $\underline S$ sont stables par $\varinjlim$ de type $\Phi$, et le foncteur d'inclusion $\underline S\longrightarrow M$ y commute, i.e. $M$ est stable pour ces limites, et $\underline S$ stable dans $M$ pour ces \unsure{$\Rightarrow$|$=$} limites.

% archive sheet number: 427

